%
% 4-homologie.tex
%
% (c) 2026 Prof Dr Andreas Müller
%
\section{Homologie
\label{buch:integration:section:homologie}}
In Anwendungen des cauchyschen Integralsatzes kommt es auf die Wahl
eines geeigneten Weges an.
Schon im Beweis des Integralsatzes wurde das Argument verwendet, dass
Wegestücke, die zweimal in entgegengesetzter Richtung durchlaufen
werden, keinen Betrag zum Integral leisten.
Ein beliebiger Weg, der einmal den Punkt $z_0$ umläuft, führt auf die
gleichen Werte des Integrals wie ein kleiner Kreis um den Punkt $z_0$.
Es scheint also gar nicht so sehr auf den Verlauf des Weges anzukommen,
sondern nur darauf was wie oft umlaufen wird.
Der Begriff der \emph{Homologie} versucht diese Beobachtung
zu präzisieren.

%
% Ketten und Zyklen
%
\subsection{Ketten und Zyklen}
Der Integralsatz von Cauchy besagt, dass ein Deformation des Weges
eines Wegintegrals sich nicht auf den Wert des Integrals auswirkt,
solange die Deformation keine Singulartität der Funktion trifft.
\input{chapters/030-integration/fig/fig-paar.tex}%
Abbildung~\ref{buch:integration:homologie:fig:paar}
zeigt die Deformation des Pfades $\gamma$, der die beiden ``Löcher''
im Definitionsbereich um die Foki $F_1$ und $F_2$ der Lemniskate
einmal umläuft, in zwei Wege $\gamma_+$ und $\gamma_-$, die $F_1$
und $F_2$ separat umlaufen, verbunden über zwei Wegstücke entlang
der reellen Achse, die in entgegengesetzter Richtung durchlaufen
werden und sich im Integral daher wegheben.
Es folgt, dass
\[
\int_\gamma f(z)\,dz
=
\int_{\gamma_+} f(z)\,dz
+
\int_{\gamma_-} f(z)\,dz.
\]
Wir müssen daher die Idee eines Weges wie $\gamma$ zu einer 
Vereinigung von Wegen wie $\{\gamma_+,\gamma_-\}$ verallgemeinern.

\begin{definition}[Kette]
\label{buch:integration:homologie:def:kette}
Eine \emph{Kette} $C$ in $U$ ist eine formale Summe
\index{Kette}%
\begin{equation}
c
=
c_1\gamma_1 + \dots + c_n\gamma_n
\label{buch:integration:homologie:eqn:def:kette}
\end{equation}
von Wegen $\gamma_k\colon I \to U$ mit $c_k\in\mathbb{Z}$.
Wir bezeichnen mit
\[
C_1
=
C_1
=
\{
c_1\gamma_1+\ldots+c_n\gamma_n
\mid
\gamma_k:I\colon U, c_k\in \mathbb{Z}
\}
\]
heisst die Menge der (eindimensionalen) Ketten.
\end{definition}

Die Definition unterscheidet alle Wege, auch solche, die ineinander
deformierbar sind.
\input{chapters/030-integration/fig/fig-kettenzyklen.tex}%
Es wird auch nicht verlangt, dass die Wege geschlossen sind.
Die Wege können geschlossen sein, es kann aber auch sein, dass sich
die Wege zu einem geschlossenen Weg zusammensetzen lassen.

\begin{definition}[Randoperator]
\label{buch:integration:homologie:definition:randoperator}
Die Menge der \emph{$0$-dimensionalen Ketten} ist
\index{Ketten!0-dimensional}%
\[
C_0 
=
C_0(U)
=
\{
c_1P_1+\ldots+c_nP_n
\mid
c_k\in \mathbb{Z}, P_k\in U
\}.
\]
Der \emph{Randoperator} ist eine
\index{Randoperator}%
lineare Abbildung $\partial \colon C_1\to C_0$, die einen
Weg $\gamma$ auf
\[
\gamma \mapsto 1\cdot \gamma(1) - 1\cdot \gamma(0)
\]
abbildet.
\end{definition}

Der Endpunkt eines Weges erhält also das Gewicht $1$, während der
Anfangspunkt das Gewicht $-1$ erhält.
Mehrere Wege können zusammengesetzt werden, wenn der Endpunkt
eines Wegestückes mit den Anfangspunkt des nächsten Stücks
übereinstimmt, wenn also $\gamma_k(1) = \gamma_{k+1}(0)$ gilt.
Der Randoperator bildet $\gamma_k + \gamma_{k+1}$ auf
\[
\partial ( \gamma_k + \gamma_{k+1} )
=
-\gamma_k(0) + \gamma_k(1) -\gamma_{k+1}(0) + \gamma_{k+1}(1)
=
-\gamma_k(0) + \gamma_{k+1}(1)
\]
ab.
Die Zusammensetzbarkeit der Wege ist also genau dann gegeben, 
wenn sich bei Anwendung des Randoperators die Punkte, wo die
Wege zusammengesetzt werden müssen, wegheben.
In Abbildung~\ref{buch:integration:homologie:fig:kettenzyklen}
tritt diese Situation im Zyklus $d$ beim Punkt $B_2$ ein.
Wenn sich alle Punkte wegheben, dann lassen sich die Wegstücke
zu einem geschlossenen Weg zusammensetzen.
In Abbildung~\ref{buch:integration:homologie:fig:kettenzyklen}
tritt dies für den Zyklus $c$ ein.
Wir definieren daher einen Zyklus wie folgt.

\begin{definition}[Zyklus]
\label{buch:integration:homologie:definition:zyklus}
Ein \emph{Zyklus} ist eine Kette in $C_1$, deren Rand verschwindet:
\index{Zyklus}%
\[
Z_1
=
Z_1(U)
=
\{
c\in C_1
\mid
\partial c = 0
\}.
\]
\end{definition}

In Abbildung~\ref{buch:integratoin:homologie:fig:kettenzyklen} sind
zwei Ketten $c=\gamma_1+\gamma_2+\gamma_3+\gamma_4$ und $d=\delta_1+\delta_2$
dargestellt.
Der Randoperator von $d$ ist
\begin{align*}
\partial d
&=
\partial\delta_1
+
\partial\delta_2
=
B_2-B_1 + B_3-B_2
=
B_3-B_1,
\intertext{es bleiben also die Puntke $B_3$ und $B_1$ mit verschiedenen
Vorzeichen.
Für $c$ ist der Randoperator}
\partial c
&=
\partial \gamma_1
+
\partial \gamma_2
+
\partial \gamma_3
+
\partial \gamma_4
=
A_2-A_1
+
A_3-A_2
+
A_4-A_3
+
A_1-A_4
=
0,
\end{align*}
die Kette $c$ ist also ein Zyklus.
Jeder Teilweg beginnt im Endpunkt eines anderen Weges und endet
im Anfangspunkt eines anderen Weges.

%
% Integral über Ketten
%
\subsection{Integral über Ketten}
Damit das Wegintegral eine lineare Abbildung von Ketten wird, 
muss das Wegintegral entlang einer Kette wird durch Linearkombination
der einzelnen Integrale gebildet:

\begin{definition}[Wegintegral einer Kette]
Ist $c\in C_1$ eine Kette der Form
\eqref{buch:integration:homologie:eqn:def:kette}
und $f$ eine in $U$ holomorphe Funktion, dann ist das 
\emph{Wegintegral entlang $C$} definiert durch
\index{Wegintegral entlang einer Kette}%
\[
\int_c f(z)\,dz
=
c_1
\int_{\gamma_1}f(z)\,dz
+
\ldots
+
c_n
\int_{\gamma_n}f(z)\,dz
=
\sum_{k=1}^n
c_k \int_{\gamma_k} f(z)\,dz.
\]
\end{definition}

Ist $F(z)$ eine Stammfunktion von $f(z)$

Ist $c$ ein Zyklus, dann lassen sich die Teilwege von $c$ zu einem
geschlossenen Weg zusammen setzen, wie dies in
Abbildung~\ref{buch:integration:homologie:fig:kettenzyklen}
dargestellt ist.
In diesem Fall ist es sinnvoll vom Integral über einen geschlossen
Weg zu sprechen und
\[
\partial c = 0
\qquad\Rightarrow\qquad
\int_c f(z)\,dz
=
\oint_c f(z)\,dz
\]
zu schreiben


Abbildung~\ref{buch:integration:homologie:fig:paar}
zeigt eine Situation einer Kette $C=\gamma_++\gamma_-$, deren
Wegintegral
\[
\int_C f(z)\,dz
=
\int_{\gamma_+} f(z)\,dz
+
\int_{\gamma_-} f(z)\,dz
=
\int_{\gamma} f(z)\,dz
\]
mit dem Wegintegral entlang der Kurve $\gamma$ übereinstimmt.
Soweit es um Wegintegrale entlang einer Kette geht, gibt es 
zwischen $\gamma$ und $C$ keinen Unterschied.

Zwischen $\gamma$ und nur der Komponente $\gamma_+$ gibt es
ein wesentlichen Unterschied.
Die Funktion $f(z) = 1/(z-F_1)$ hat eine Polstelle im Punkt $F_1$.
Die Wegintegrale sind
\begin{align*}
\frac{1}{2\pi i}
\oint_\gamma f(z)\,dz
&=
\\
\frac{1}{2\pi i}
\oint_\gamma \frac{1}{z-F_1}\,dz
&=
\\
\operatorname{ind}_\gamma(F_1)
&=
1
\\
&\ne
\frac{1}{2\pi i}
\oint_{\gamma_+} f(z)\,dz
\\
&=
\frac{1}{2\pi i}
\oint_{\gamma_+} \frac{1}{z-F_1}\,dz
\\
&=
\operatorname{ind}_{\gamma_+}(F_1)
\\
&=
0
\end{align*}
sind jedoch verschieden.
Der Grund dafür ist, dass $\gamma$ und $\gamma_+$ den Punkt $F_1$
nicht gleich oft umlaufen.
Wir definieren daher ein Äquivalenzrelation von Ketten durch die
Anzahl der Male, die sie Punkte ausserhalb von $U$ umlaufen.

\begin{definition}[Umlaufzahl einer Kette]
Die \emph{Umlaufzahl} einer Kette $C$ um einen Punkt $z_0$ ist
\index{Umlaufzahl}%
\[
\operatorname{ind}_C(z_0)
=
\frac{1}{2\pi i}
\int_C
\frac{1}{z-z_0}
\,dz.
\]
\end{definition}

%
% Homologie von Ketten
%
\subsection{Homologie von Ketten
\label{buch:integration:homologie:subsection:homologie}}
Die Homotpieeigenschaft des Wegintegrals hat gezeigt, dass sich der
Wert eines Wegintegrals einer holomorphen Funktion nicht ändert, wenn
man den Weg stetig deformiert.
Dies wurde zum Beispiel in Abbildung~\ref{buch:integration:fig:residuum}
verwendet, um zu zeigen, dass der Weg $\gamma$ und der kleine Kreis
um den Punkt $z_0$ zu den gleichen Werten der Wegintegrale führen.
Diese Äquivalenz von Ketten, die zu gleichen Wegintegralen führen,
verdient einen eigenen Namen.

\begin{definition}[homolog]
Die Ketten $c_1$ und $c_2$ heissen \emph{homolog} in $U$, in Zeichen
$c_1\sim c_2$ oder $c_1\sim_{U}c_2$, wenn
\index{homolog}%
für jeden Punkt $a\in \mathbb{C}\setminus U$
\[
\operatorname{ind}_{c_1}(a)
=
\operatorname{ind}_{c_2}(a)
\]
gilt.
\end{definition}

Die Kette $\gamma_++\gamma_-$ ist also zu $\gamma$ homolog.
Für homologe Ketten ist nicht nur die Umlaufzahl gleich, sondern
jedes Integral.

\begin{satz}
Sind die Ketten $c_1$ und $c_2$ homolog, dann ist
\[
\int_{c_1} f(z)\,dz
=
\int_{c_2} f(z)\,dz.
\]
\end{satz}

Wir betrachten jetzt einen später besonders nützlichen Spezialfall.
\input{chapters/030-integration/fig/fig-merohomolog.tex}%
In Abbildung~\ref{buch:integration:homologie:fig:merohomolog}
ist ein Gebiet $U$ mit einer einfach geschlossenen Kurve
$\gamma$ dargestellt.
Das Komplement von $U$ im Inneren der Kurve $\gamma$ besteht
aus den isolierten Punkten $z_k$.
Jeder dieser Punkte kann so mit der Kurve $\gamma$ verbunden werden,
dass sich die Verbindungen nicht schneiden.
Ausserdem wird jeder von diesen Punkten von einem kleinen Kreis $\gamma_k$
im Gegenuhrzeigersinn umlaufen.
Aus $\gamma$, den im Uhrzeigersinn durchlaufenen Kreisen $-\gamma_k$
und den Verbindungen zu $\gamma$ lässt sich ein neuer Weg zusammensetzen,
der sich in $U$ zusammenziehen lässt.
Das Wegintegral über den gesamten Weg verschwindet daher.
Die Verbindungswege zur Kurve $\gamma$ werden dabei zweimal in
entgegengesetzter Richtung durchlaufen und tragen daher nicht zum
Wegintegral bei.
Das Wegintegral einer in $U$ holomorphen Funktion über $\gamma$
ist daher gleich der Summe der Wegintegrale der Funktion
über die Kreise $\gamma_k$.

\begin{satz}
\label{buch:integration:homologie:satz:pfadpunkte}
Ist $\gamma$ eine einfach geschlossene Kurve in $U$ und besteht das Komplement
von $U$ im Inneren von $\gamma$ nur aus endlich vielen isolierten Punkten
$z_k$, $1\le k\le n$, dann gibt es einen Radius $r$ derart, dass $\gamma$
homolog zu einem Zyklus 
\[
c
=
\sum_{k=1}^n \gamma_k,
\]
wobei jedes $\gamma_k$ ein Kreis vom Radius $r$ um $z_k$ ist.
\end{satz}

\begin{proof}
Es muss nur der Radius $r$ berechnet werden.
Nimmt man
\[
r
=
\frac13
\min \{ |z_k-z_l|\mid 1\le k,l\le n\},
\]
dann können sich die Kreise um die Punkte $z_k$ mit Radius $r$ nicht
schneiden.
Für jeden Punkt $z_k$ ist dann die Umlaufzahl 
\[
\operatorname{ind}_{\gamma_k}(z_l)
=
\delta_{kl}
=
\begin{cases}
1&\qquad\text{für $k=l$}\\
0&\qquad\text{für $k\ne l$.}
\end{cases}
\]
Für den ganzen Zyklus $C$ gilt dann
\[
\operatorname{ind}_c(z_l)
=
\sum_{k=1}^n
\delta_{kl}
=
1
=
\operatorname{ind}_\gamma(z_l).
\]
Dies zeigt, dass $\gamma$ und $c$ homolog sind.
\end{proof}

%
% Homologiegruppe
%
\subsection{Homologiegruppe}
Im Abschnitt \ref{buch:integration:homologie:subsection:homologie}
haben wir die Homologierelation durch die die Umlaufzahl um Punkte,
die ausserhalb des Definitionsgebietes liegen, definiert.
\input{chapters/030-integration/fig/fig-raender.tex}%
Unter Verwendung der Homotopieeigenschaft können wir aber auch eine
algebraische Definition aufbauen.
Dazu konstruieren wir zunächst die zweidimensionalen Ketten.

\begin{definition}[Zweidimensionale Ketten]
Eine \emph{zweidimensionale Kette} ist eine formale Linearkombination
von Abbildungen 
\[
\Gamma
\colon
I\times I \to U
:
(s,t) \mapsto \Gamma(s,t)
\].
Der \emph{Randoperator} für zweidimensionale Ketten ist die lineare
Abbildung, die einer einzelnen Abbildung $\Gamma$ die Summe
\[
\partial \Gamma
=
\gamma_1
+
\gamma_2
+
\gamma_3
+
\gamma_4
\]
zuordnet, wobei die $\gamma_k$ die Wege
\[
\begin{array}{rlcl}
\gamma_1 &\colon I \to U&:&t\mapsto \Gamma(t,0)\\
\gamma_2 &\colon I \to U&:&t\mapsto \Gamma(1,t)\\
\gamma_3 &\colon I \to U&:&t\mapsto \Gamma(1-t,1)\\
\gamma_4 &\colon I \to U&:&t\mapsto \Gamma(0,1-t)
\end{array}
\]
sind.
\end{definition}

Die Wege $\gamma_k$ sind in
Abbildung~\ref{buch:integration:homologie:fig:raender} dargestellt.
$\gamma_1$ ist der Weg auf dem Rand von $A$ nach $B$,
$\gamma_2$ ist das Randstück von $B$ nach $C$,
$\gamma_3$ das Randstück von $C$ nach $D$ und
$\gamma_4$ das Randstück von $D$ zurück nach $A$.
Es ist offensichtlich, dass die vier Wegstücke einen geschlossenen
Weg bilden und dass daher $\partial\partial \Gamma=0$ ist.
Die Ketten $\gamma_1$ und $-\gamma_2-\gamma_3-\gamma_4$ sind homolog
da die $\gamma_1+\gamma_2+\gamma_3+\gamma_4$ ein geschlossener
Weg ist.


%
% Homologie in höheren Dimensionen
%
\subsection{Homologie in höheren Dimensionen}



